\section{Consequences and directions}\label{sec:directions}

\subsection{What follows at once}
The truth of \cref{st:rh,st:grh,st:strip} has many known consequences, and none of them depends
on how the theorem was proved. We list those that seem most useful to state formally; the
first three are proved in a short Lean file published with this paper, against the compiled
archive, with the same three axioms.
\begin{itemize}
\item $\zeta(s)\ne0$ for $\re s>\tfrac12$, $s\neq1$.
\item No exceptional zeros: for a non-trivial character $\chi$, $L(\sigma,\chi)\neq0$ for real
  $\sigma>\tfrac12$.
\item Zeros of $L(s,\chi)$ in the open critical strip lie on the critical line.
\end{itemize}
The following are stated there as targets and are not proved by us.
\begin{itemize}
\item The prime number theorem in the form $\pi(x)=\mathrm{li}(x)+O(\sqrt x\log x)$
  \cite{vonkoch1901}, and primes in intervals $(x,x+C\sqrt x\log x]$ \cite{cramer1920}.
\item Miller's deterministic primality test \cite{miller1976} with the explicit bound of Bach
  \cite{bach1990}: an odd $n>2$ is prime if and only if it is a strong probable prime to every
  base $a$ with $2\le a<n$ and $a\le2(\log n)^2$. This becomes unconditional. It is slower in practice than the
  probabilistic tests in use; beside the unconditional test of Agrawal, Kayal and Saxena
  \cite{aks2004} its merit is simplicity. The gain is a short deterministic test whose
  correctness can now be proved formally.
\item Bounds for the least quadratic non-residue \cite{ankeny1952} and the least prime in an
  arithmetic progression \cite{bachsorenson1996}.
\end{itemize}

\subsection{What does not follow}
No method of factoring integers follows, and nothing in the argument suggests one. The
statement was already assumed by everyone who would use it, so its truth changes
``fast if GRH holds'' into ``fast'' and makes nothing faster than was believed. The argument
itself bounds averaged counts of primes by means of auxiliary fields of degree about $x^2$ and
conductor exponential in a power of $x$; it never computes with a single integer.

The scope should also be kept in view. The theorem covers Dirichlet $L$-functions. Results
that need the Riemann Hypothesis for Dedekind zeta functions of non-abelian fields, for
elliptic curves or for Hecke characters are not covered.

\subsection{Directions}
\emph{An intelligible proof.} The present account describes the proof; it does not make it
short. A statement of the whole argument on one page, each step a statement of the archive, is
published with this paper (\lean{one-page.pdf} in the repository of \cref{app:reproduce}). It
is a skeleton, not a proof that can be checked by hand. Two statements seem to us to be the
right core of a human proof: the slope inequality
for nearly tight frames (\cref{st:F}), which is self-contained, and the identity
(\cref{st:key}) that turns the period side into the analytic main term. A proof organised
around them, with the divisibility of \cref{sec:divisibility} as its one arithmetic input and
everything else quoted from the literature, would be the next thing to write.

\emph{What is not needed.} The dependency cone shows that about a third of the archive's
theorems are not used by the final theorem (\cref{sec:offpath}), and several inputs are used far below their
strength: the cohomological bound is needed only in the form ``at most $\exp(x^\eps)$''
(\cref{sec:cohomology}), and the term $\log\abs I$ in \cref{st:F} is slack. A systematic
weakening of each input to what is used would shorten the proof and show which of the deep
theorems are essential.

\emph{The by-products.} If the archive's proofs of the published theorems it relies on are what
their names say, it contains formalisations of a substantial part of modern number theory:
class field theory, Galois representations for unitary groups, the Siegel--Weil formula,
arithmetic intersection theory, resolution of singularities. Each deserves its own audit of
statement against literature, and extraction into a form that others can use.

\emph{The cohomological bound.} \Cref{st:vwin} is a new kind of result about the groups that
Vandiver's conjecture concerns. Whether its method reaches the plus part of the class groups of
cyclotomic fields themselves, which some arguments in lattice-based cryptography assume to be
small \cite{cdpr2016}, we do not know.

\emph{Other $L$-functions.} The note says that Dedekind zeta functions are expected to yield to
the same method.
